\section{レンズ(1)}

\subsection{凸レンズ}
\begin{wrapfigure}[13]{r}[5mm]{70mm}
  \centering
  \vspace{-5mm}
  \includegraphics[keepaspectratio,width=70mm]
  {text2_p15_fig1.png}
\end{wrapfigure}
レンズは光の屈折を利用して光を集めたり発散させることができるもので，2つの球面によって挟まれたガラスなどでできている．
レンズには凸レンズと凹レンズがあるが，まず凸レンズについて考えてみよう．

レンズの2つの球面の中心を結ぶ直線をレンズの{\bf 光軸}という．
光軸に対して平行な光線を凸レンズにあてると，光はレンズの後方（向こう側）のある一点Fに集まる．
この点Fをレンズの{\bf 焦点}といい，レンズの中心から焦点までの距離をレンズの{\bf 焦点距離}という．
レンズの焦点はレンズの前方と後方に一つずつあり，いずれの焦点距離も等しい．
焦点に光源を置いた場合は先ほどのちょうど逆となり，レンズを通ると光軸に平行な光線群となる．\\[10mm]


次に，レンズより小さいサイズの物体を光軸上，手前の焦点よりも外側に置き，物体のある点Pから出ていく以下の3つの光について考える．\vspace{5mm}

\begin{wrapfigure}[8]{r}[5mm]{90mm}
  \centering
  \vspace{-5mm}
  \includegraphics[keepaspectratio,width=90mm]
  {text2_p15_fig2.png}
\end{wrapfigure}
$\MARU{1}$\hspace{2mm} {\bf 光軸に平行に出た光}

この光はレンズの向こう側の焦点を通るように進む．

$\MARU{2}$\hspace{2mm} {\bf レンズの中心を通過する光}

この光は曲がらずに直進する．（レンズの厚さが無視できるとき）

$\MARU{3}$\hspace{2mm} {\bf 手前側の焦点を通る光}

この光はレンズの向こう側に出ると光軸に平行な光線となる．\vspace{5mm}

今は3つのみについて考えたが，点Pから出てレンズに入った光は$\MARU{1},\MARU{2},\MARU{3}$に限らず，{\bf すべて同一の点Qを通る．}
そのため，光軸に対して垂直でQ点を含む面に紙を置くと，紙面には光源の像が生じる．
この像のように，実際に光が集まって作られる像を{\bf 実像}と呼ぶ．
この場合，像は光源の物体に対して逆さまになっているので，これを{\bf 倒立像}と呼ぶ．
これに対し像が同じ向きになっているものは{\bf 正立像}と呼ぶ．

なお，実際の像の位置の把握には，上述の3つのうち2つを考えれば十分である．

\newpage
\subsection{凹レンズ}
\begin{wrapfigure}[15]{r}[5mm]{70mm}
  \centering
  \vspace{-5mm}
  \includegraphics[keepaspectratio,width=70mm]
  {text2_p16_fig1.png}
\end{wrapfigure}
凹レンズは周辺部よりも中心部が凹んでいる形のレンズである．
凹レンズの光軸に対して平行な光線をあてると，光はレンズを通ると図のように広がっていく．
この広がっていく光線をレンズの手前側まで伸ばすと，直線はある一点で交わる．
すなわち，レンズにより発散した光線はこの一点から直進して出ていった光と同様になる．
この点を凹レンズの{\bf 焦点}という．

凹レンズについても凸レンズ同様，レンズの前方と後方に焦点が一つずつあり，その焦点距離はいずれも等しい．
またレンズの向こう側にある焦点に光を集めるように凹レンズに光を入射すると，その光は凹レンズを通過した後，光軸に平行な光線群となる．\\[10mm]

次にレンズより小さいサイズの物体を光軸上，焦点よりも外側に置き，物体のある点Pから出ていく以下の3つの光について考える．\vspace{5mm}

\begin{wrapfigure}[8]{r}[5mm]{90mm}
  \centering
  \vspace{-5mm}
  \includegraphics[keepaspectratio,width=90mm]
  {text2_p17_fig1.png}
\end{wrapfigure}
$\MARU{1}$\hspace{2mm} {\bf 光軸に平行に出た光}

この光はレンズの手前側の焦点から直進するように進む．

$\MARU{2}$\hspace{2mm} {\bf レンズの中心を通過する光}

この光は曲がらずに直進する．（レンズの厚さが無視できるとき）

$\MARU{3}$\hspace{2mm} {\bf 向こう側の焦点へ直進する光}

この光はレンズの向こう側に出ると光軸に平行な光線となる．\vspace{5mm}

レンズの向こう側に出た$\MARU{1} 〜\MARU{3}$の3つの光線はレンズの向こう側では交わらないが，これらの光線をレンズの手前側に戻してやると，それらは図のQ点の一点に交わる．
すなわち，レンズの向こう側の光線は，Qの位置から直進して出てきた光のように広がっていく．
そのため，レンズの向こう側からレンズを覗き見るとあたかもQに物体があるかのように見える．
このQを{\bf 虚像}という．
虚像は実像と異なり，実際にその位置に光が集まってはいないため虚像の位置に紙を置いても像は写らない．

\newpage
\subsection{虚光源}
\begin{wrapfigure}[10]{r}[5mm]{70mm}
  \centering
  \vspace{-5mm}
  \includegraphics[keepaspectratio,width=70mm]
  {text2_p19_fig1.png}
\end{wrapfigure}
凹レンズに対し，レンズの向こう側にある焦点に光を集めるように凹レンズの光を入射すると，その光は凹レンズを通過した後，光軸に平行な光線群となる．

ここで考えた，レンズの向こう側にある焦点に向かう光は，凹レンズよりも手前に凸レンズを置くことにより作ることができる．

このように，レンズの向こう側にある一点に向かう光線群において，集光しようとしている場所を{\bf 虚光源}と呼ぶ．

光源，像の位置を次のように定義し直すと，虚光源や虚像もまとめて定義することができる．

   \begin{matome}{光源と像}
   
\hspace{2mm}$\MARU{1}$\hspace{2mm} {\bf 光源}\\
 \hspace{3mm} {\bf レンズに入射する直前の光線群（またはその延長）が一点に集まる場所．}
 光の進行方向に対して，その場所がレンズより手前ならば{\bf 実光源}，レンズより奥ならば{\bf 虚光源}となる．
 実際に光線群が光源の位置で集光しているかどうかは光源の位置を決定する上で関係がない．

\hspace{2mm}$\MARU{2}$\hspace{2mm} {\bf 像}\\
 \hspace{3mm} {\bf レンズにより屈折された直後の光線群（またはその延長）が一点に集まる場所．}
光の進行方向に対して，その場所がレンズよりも奥ならば{\bf 実像}，レンズより手前ならば{\bf 虚像}となる．
 実際に光線群が像の位置で集光しているかどうかは像の位置を決定する上で関係がない．\\[27mm]
 \end{matome}
 \begin{figure}[h]
\begin{center}
\vspace{-40mm}
\includegraphics[width=130mm]{text2_p19_fig2.png}
\end{center}
\end{figure}

 レンズを通じて，入射する光が曲げられ，集光する位置が変化することを{\bf 写像}と呼ぶ．
 光源，像は1回の写像ごとに位置が定義されるもので，特に複数のレンズを組み合わせた光学系では，各写像に対し，入射する光が実際に存在する光源から出た光であるとは限らず，また，レンズを通過した光が一点に集光するとは限らない．
 
 \newpage
 \subsection{凹面鏡，凸面鏡}
 
 鏡によっても，光の集光，写像を五輪することができる．
 中央部がへこんでいる鏡を凹面鏡，出っ張っている鏡を凸面鏡と呼ぶ．
 
 凹面鏡に光軸に対して平行な光線をあてると，反射光は鏡の手前のある一点に集光する．
 この点を凹面鏡の焦点と呼び，凹面鏡から焦点までの距離を焦点距離と呼ぶ．
 また，焦点の位置に光源を置くと，光線の通り道は先ほどと逆になり，反射光は光軸と平行な光線群となる．
 
 凸面鏡に光軸に対して平行な光線をあてると，反射光は広がっていき，その光線を凸面鏡の奥側まで伸ばすと，直線は一点で交わる．
 この位置を凸面鏡の焦点と呼び，凸面鏡から焦点までの距離を焦点距離と呼ぶ．
 また，凸面鏡の奥側の焦点に集まるような光線を入射すると，反射光は光軸と平行な光線群となる． 
  \begin{figure}[h]
\begin{center}
\includegraphics[width=130mm]{text2_p21_fig1.png}
\end{center}
\end{figure}


\newpage

\begin{figure}[h]
\begin{center}
\includegraphics[width=170mm]{prob2_p11.png}
\vspace{150mm}
\end{center}
\end{figure}

\newpage

\begin{figure}[h]
\begin{center}
\includegraphics[width=170mm]{prob2_p12.png}
\vspace{150mm}
\end{center}
\end{figure}

\newpage

\begin{figure}[h]
\begin{center}
\includegraphics[width=170mm]{prob2_p13.png}
\vspace{150mm}
\end{center}
\end{figure}

\newpage

\begin{figure}[h]
\begin{center}
\includegraphics[width=170mm]{prob2_p83_2.png}
\vspace{150mm}
\end{center}
\end{figure}

\newpage

\begin{figure}[h]
\begin{center}
\includegraphics[width=170mm]{prob2_p84.png}
\vspace{150mm}
\end{center}
\end{figure}

\newpage

\begin{figure}[h]
\begin{center}
\includegraphics[width=170mm]{prob2_p85.png}
\vspace{150mm}
\end{center}
\end{figure}

\newpage

\begin{figure}[h]
\begin{center}
\includegraphics[width=170mm]{prob2_p86.png}
\vspace{150mm}
\end{center}
\end{figure}
